\honest{A Fable of Science and Politics}{Eliezer Yudkowsky}{2006}

In the Roman Empire, the Blues and the Greens murdered each other in duels, ambushes, battles and riots. Procopius wrote that their hostility ``has no cause'' and survives marriage, kinship and friendship. Gibbon wrote that every candidate for office needed the support of a faction. Who were they? They were ``sports fans,'' supporters of two chariot-racing teams. \nb{Procopius, in the passage quoted, says they fought over the names and the seats at the games. But historians disagree about how far the factions had become political and religious parties by then, and Gibbon's next sentence, which the post omits, links the Greens to the sect of Anastasius and the Blues to orthodoxy.}

Now imagine a society that has sealed itself underground. An old scrap of paper says the sky is ``cerulean,'' and the society splits over whether that means blue or green. After long wars there is a truce. The Blues and Greens still take opposite sides on taxes, marriage laws and the shape of the Earth, and the positions come in packages: it would be rare to find a city merchant who thought the sky blue and yet backed the Green taxes. But they shop in each other's shops, and tolerance is growing. \nb{Taxes and marriage laws are questions of value. The discovery that follows can settle only the questions of fact.}

An earthquake opens an old passage, and one member of a sightseeing party of six climbs to the surface. The sky is blue. The story then branches six ways, one for each person who might have climbed.

Aditya, a Blue, smiles with hatred and goes back to end the truce. Barron, a Green, remembers the massacres of Greens and weeps, sure that the universe ``had always been a place of evil.'' Charles, a Blue professor who taught that both views were ``equally valid,'' wonders for a moment whether a Green would see a green sky here, then decides to block the passage, fearing the sky will be ``misinterpreted.'' \nb{Aditya's branch shows the danger Charles fears. The post does not say whether hiding the sky is wrong.} Daria, a Green, remembers her father's advice to speak the truth even if her voice trembles, and forces herself to accept that she was wrong, knowing her family may turn from her: ``it's not so complicated, after all.'' Eddin, a Green, laughs at a history of murder over ``one pointless thing,'' and wonders how to keep the news from blowing up the world, and whether they all deserve it. Ferris, whose faction the story does not give, is struck with ``sheer wonder and delight,'' says ``Oh, so that's what color it is,'' and goes exploring.

I state no moral. \nb{In the book, ``The Meditation on Curiosity'' says Ferris ``embodies the power of innocent curiosity.'' The post gives that response only to the one character whose side it does not name. Daria is the only partisan who reaches the truth, and she does it in grief. The post does not say how a partisan could become Ferris.}
